\documentclass[12pt]{article}
\usepackage[margin=0.82in]{geometry}
\usepackage[T1]{fontenc}
\usepackage[utf8]{inputenc}
\usepackage{amsmath,amssymb}
\usepackage{booktabs,array}
\usepackage{xcolor}
\usepackage{tikz}
\usepackage{enumitem}
\usepackage{fancyhdr}
\usepackage{mdframed}
\definecolor{navy}{RGB}{38,69,96}
\definecolor{pale}{RGB}{235,242,247}
\definecolor{gold}{RGB}{166,126,30}
\pagestyle{fancy}\fancyhf{}
\lhead{MATH\& 107 \quad Module 5}
\rhead{Cybernetic Systems}
\cfoot{\thepage}
\setlength{\headheight}{15pt}
\setlist[enumerate]{itemsep=8pt,topsep=5pt}
\newcommand{\answerline}{\par\vspace{0.38in}\hrule\vspace{0.12in}}
\newcommand{\shortline}{\par\vspace{0.26in}\hrule\vspace{0.1in}}
\begin{document}

\begin{center}
{\Large\bfseries\color{navy}Module 5: Information, Feedback, and Cybernetic Systems}\\[4pt]
{\large Student Summary and Practice}
\end{center}
\vspace{2mm}
\noindent Name: \rule{3.1in}{0.4pt}\hfill Date: \rule{1.4in}{0.4pt}

\begin{mdframed}[backgroundcolor=pale,linecolor=navy]
\textbf{The question for this module.} How can a system use information about what just happened to guide what happens next? A thermostat, a driver, a greenhouse controller, and a learning algorithm can all be studied with the same pattern:
\[
\boxed{\text{state}\;\longrightarrow\;\text{measurement}\;\longrightarrow\;\text{comparison}\;\longrightarrow\;\text{action}\;\longrightarrow\;\text{new state}}.
\]
The new state can be measured again, closing the loop. A system organized this way is called a \textbf{cybernetic system}.
\end{mdframed}

\section*{1. Read the signal}
The \textbf{state} is the condition that can change. A \textbf{measurement} is the information the system receives about that state. A \textbf{target} or \textbf{reference value} is the condition the system is trying to reach. The difference between the target and the measured value is the \textbf{error signal}:
\[
e_n=r-x_n,
\]
where $r$ is the target and $x_n$ is the measured state at step $n$. Here, ``error'' means a difference; it does not necessarily mean a mistake.

\textbf{Example.} A classroom is at $64^\circ$F and the thermostat target is $70^\circ$F. Then $e=70-64=6^\circ$F. The room is six degrees below its target, so the controller calls for heat.

\section*{2. Corrective feedback can be gentle or too strong}
A simple correction rule is
\[
x_{n+1}=x_n+k(r-x_n),
\]
where $k$ is the response strength. It tells the system what fraction or multiple of the current error to correct.

\textbf{Example.} With target $r=70$, starting state $x_0=64$, and $k=0.5$, the system corrects half of the remaining error each step:
\[
\begin{array}{c|c|c|c}
\toprule
\text{Step }n & \text{State }x_n & \text{Error }70-x_n & \text{Next state }x_{n+1}\\
\midrule
0&64&6&64+0.5(6)=67\\
1&67&3&67+0.5(3)=68.5\\
2&68.5&1.5&68.5+0.5(1.5)=69.25\\
\bottomrule
\end{array}
\]
For this model, $0<k<1$ approaches the target without crossing it. If $1<k<2$, the system crosses the target and then corrects back; the overshoots shrink. At $k>2$, the overshoots grow. \textbf{Stability} describes whether departures settle down or grow.

\section*{3. Feedback can amplify change}
\textbf{Corrective (negative) feedback} acts against a departure from a target. The word negative describes the direction of the response, not whether it is harmful. \textbf{Amplifying (positive) feedback} reinforces a change, so a change can produce more change in the same direction.

For a simple percentage-growth process,
\[
x_{n+1}=(1+g)x_n,
\]
where $g$ is the growth rate. Starting at $100$ with $g=0.20$ gives $100,120,144,\ldots$. Each increase is calculated from the larger value that came before it.

\section*{4. Information is what the system can use}
A controller acts on a \textbf{signal}, not directly on the world. A signal can be rounded, noisy, incomplete, or delayed. If the controller sees an old measurement, it may correct a condition that has already changed and overshoot.

Information can be measured in bits. For an event with probability $p$, its surprisal is
\[
I=-\log_2(p).
\]
An event with probability $1/8$ has $I=-\log_2(1/8)=3$ bits. A parity bit can help detect some transmission errors, but it does not tell the receiver which bit changed or repair the message. These tools matter because a feedback loop is only as useful as the information it receives.

\section*{5. Learning is a feedback process}
A prediction by itself is a one-way calculation. Learning closes a loop: make a prediction, compare it with a target, measure the error, change adjustable weights, and predict again. In a real neural network, backpropagation calculates how the final error is connected to weights in earlier layers. A training rule uses that information to update the weights.

\begin{mdframed}[backgroundcolor=pale,linecolor=gold]
\centering\textbf{attempt}\ $\longrightarrow$ \textbf{error}\ $\longrightarrow$ \textbf{change}\ $\longrightarrow$ \textbf{new attempt}
\end{mdframed}

\section*{Quick check}
\begin{tabular}{>{\raggedright\arraybackslash}p{0.23\linewidth}p{0.68\linewidth}}
\toprule
\textbf{Idea} & \textbf{Ask yourself}\\
\midrule
Feedback loop & What is measured, compared, and changed?\\
Error & Is the measurement above or below the target?\\
Response strength & Does the correction approach, cross, or move away from the target?\\
Signal quality & Is the information current and precise enough for the action?\\
Learning & Does the error from one attempt change the next attempt?\\
\bottomrule
\end{tabular}

\newpage
\begin{center}
{\Large\bfseries\color{navy}Module 5 Practice}
\end{center}
\noindent Show your calculations. Then explain what each number means in the system.

\begin{enumerate}
\item A greenhouse controller tries to keep the room at $72^\circ$F. Its sensor reports $66^\circ$F. Identify the target, the measurement, and the error. State what the positive error means.
\answerline
\item The target is $100$, the current state is $80$, and the controller uses $k=0.25$.
  \begin{enumerate}[label=(\alph*)]
  \item Find the error.
  \item Find the next state using $x_{n+1}=x_n+k(r-x_n)$.
  \item How much closer to the target did the system move?
  \end{enumerate}
\answerline
\item A controller uses the same rule with $k=1.5$. The target is $100$ and the current state is $80$.
  \begin{enumerate}[label=(\alph*)]
  \item Find the next state.
  \item Did the system cross the target?
  \item In this simple model, what happens to later overshoots when $1<k<2$?
  \end{enumerate}
\answerline
\item A store wants to keep $100$ items in stock. Its current inventory is $60$, but the controller acts on yesterday's report of $80$ and orders half the reported shortage.
  \begin{enumerate}[label=(\alph*)]
  \item How many items does the controller order from the delayed report?
  \item How many would it order if it used the current inventory instead?
  \item Explain why delayed information can lead to a poor correction.
  \end{enumerate}
\answerline
\item A one-weight learning model predicts $\widehat y=2w$. The target is $6$, and the update rule is $w_{\text{new}}=w+0.25(6-\widehat y)$. Begin with $w=1$.
  \begin{enumerate}[label=(\alph*)]
  \item Find the prediction and error on the first attempt.
  \item Find the updated weight.
  \item Find the next prediction. Did it move closer to the target?
  \end{enumerate}
\end{enumerate}

\end{document}
